\chapter{Calcium (Urine)}

\section*{What Is This Test?}
The urine calcium test measures how much calcium is excreted in the urine. The standard method uses a 24-hour urine collection (24-hour urine calcium). Less commonly, a random spot sample measures the calcium-to-creatinine ratio. Urine calcium reflects the net result of four processes. These are intestinal calcium absorption, bone turnover, renal calcium filtration, and renal tubular calcium reabsorption. Renal calcium handling is regulated mainly by PTH and calcitriol. PTH increases distal tubular reabsorption. Calcitriol increases intestinal absorption. Urine calcium measurement is used for three main purposes. First, it evaluates kidney stone disease (nephrolithiasis) to identify hypercalciuria and guide preventive therapy. Second, it differentiates familial hypocalciuric hypercalcemia (FHH) from primary hyperparathyroidism in patients with hypercalcemia. Third, it monitors patients with nephrocalcinosis. It also monitors patients on treatments that affect renal calcium handling (loop diuretics, thiazides, vitamin D). Urine calcium is not a measure of serum calcium. It does not directly reflect calcium intake. It is primarily a measure of renal calcium excretion. A 24-hour urine collection also provides urine volume, sodium, citrate, oxalate, uric acid, and creatinine. These are essential for the complete metabolic evaluation of kidney stones.

\section*{Alternative Names}
\begin{itemize}
  \item 24-Hour Urine Calcium
  \item Urinary Calcium
  \item Urine Calcium/Creatinine Ratio
  \item Calcium Excretion (Urine)
  \item 24-Hour Urine Stone Risk Panel
\end{itemize}
\section*{Principle}
Urine calcium is measured by the same colorimetric methods used for serum calcium (Arsenazo III or o-cresolphthalein complexone). It can also be measured by ion-selective electrodes. It is reported in mg/24 hours. The 24-hour urine collection is the standard. The patient voids at the start time and discards the first void. The patient then collects all urine for 24 hours, including the first void at the end time. The container is kept refrigerated. The specimen is delivered to the laboratory. The calcium-to-creatinine ratio (mg calcium/mg creatinine) can be measured on a random spot sample. It is useful when a 24-hour collection is impractical, especially in children. The fractional excretion of calcium (FECa) is calculated as: FECa = (urine Ca $\times$ serum Cr) / (serum Ca $\times$ urine Cr) $\times$ 100\%. In familial hypocalciuric hypercalcemia, the calcium-to-creatinine clearance ratio (CCCR) is <0.01. In primary hyperparathyroidism, it is typically >0.01 (usually >0.02). Hypercalciuria is defined as >300 mg/24 h in men and >250 mg/24 h in women (some laboratories use >4 mg/kg/day). Urine calcium excretion is influenced by several factors. Dietary sodium and protein both increase calcium excretion. Calcium intake and acid-base status also influence urine calcium.

\section*{History}
The association between hypercalciuria and kidney stone disease was recognized in the mid-20th century. Pak and colleagues (1979) classified the mechanisms of hypercalciuria into absorptive, renal, and resorptive subtypes \cite{pak1979}. They established the concept of a selective 24-hour urine evaluation for stone formers. The distinction between familial hypocalciuric hypercalcemia (FHH) and primary hyperparathyroidism was recognized in the 1970s. The 24-hour urine calcium (or calcium-to-creatinine clearance ratio) became the standard discriminator. Modern guidelines for the metabolic evaluation of nephrolithiasis recommend a 24-hour urine collection \cite{worcester2010}\cite{pearle2014}. The panel includes calcium, oxalate, citrate, uric acid, sodium, creatinine, and volume. It is recommended for recurrent or high-risk stone formers.

\section*{Why This Test Is Ordered}
\begin{itemize}
  \item \textbf{Evaluation of kidney stones (nephrolithiasis):} 24-hour urine calcium is part of the metabolic evaluation of calcium oxalate and calcium phosphate stones. It identifies hypercalciuria, the most common metabolic abnormality in stone formers (present in 30--60\%). It guides preventive therapy. Thiazides reduce urinary calcium excretion by 40--50\% and reduce stone recurrence.
  \item \textbf{Differentiating familial hypocalciuric hypercalcemia (FHH) from primary hyperparathyroidism:} In a patient with hypercalcemia, a low 24-hour urine calcium (or low calcium-to-creatinine clearance ratio <0.01) supports FHH. Normal or high urine calcium supports primary hyperparathyroidism.
  \item \textbf{Monitoring treatment:} Urine calcium is monitored in patients on thiazide diuretics for hypercalciuria. It is also monitored in patients receiving vitamin D or calcium supplementation to detect hypercalciuria.
  \item \textbf{Evaluation of nephrocalcinosis:} Renal calcium deposition detected on imaging warrants a metabolic workup including urine calcium.
  \item \textbf{Complete metabolic evaluation of high-risk stone formers:} This applies to recurrent stones, bilateral stones, stones in children, family history of stones, or single stones in a high-risk individual (for example, a pediatric or solitary-kidney patient). The 24-hour urine panel includes calcium, oxalate, citrate, uric acid, sodium, and volume.
\end{itemize}

\section*{Primary vs Secondary Test}
Urine calcium is a primary test for the metabolic evaluation of nephrolithiasis. It is also a primary test for distinguishing FHH from primary hyperparathyroidism. It is a secondary/supplementary test in the routine evaluation of hypercalcemia.

\section*{How This Test Is Performed}
For the 24-hour urine test, the patient collects all urine over 24 hours in a large collection jug. The jug is typically 2.5--3 liters and is provided by the laboratory. The container may contain a small amount of preservative (usually hydrochloric acid) to keep calcium dissolved. The collection starts by discarding the first morning void. The patient then collects all urine, including the final morning void of the next day. The container is refrigerated during collection. It is transported to the laboratory promptly. For a spot calcium-to-creatinine ratio, a single random urine sample is collected, usually in the fasting state.

\section*{What Preparation Is Needed}
Patients should not take calcium or vitamin D supplements for 1 week before the 24-hour collection, if possible. This reflects baseline excretion. They should maintain their usual diet, including usual calcium intake (about 1000 mg/day). They should avoid excessive dairy intake. They should drink their usual amount of water. Medications that affect calcium excretion (thiazides, loop diuretics, corticosteroids, bisphosphonates) should be documented. The provider decides whether to withhold them. Complete 24-hour collection is essential. The result is invalid if any urine is missed or spilled. Laboratory verification of collection adequacy uses total urine creatinine. Men typically excrete 15--25 mg/kg/day and women 10--20 mg/kg/day.

\section*{Contraindications and Precautions}
\begin{itemize}
  \item Incomplete 24-hour urine collection invalidates the result. The most common error is missing one or more voids, which falsely lowers calcium. Adequacy is checked by total urine creatinine.
  \item Urine calcium is affected by diet. High sodium intake increases calcium excretion (each 100 mEq sodium increases urine calcium by approximately 20--40 mg/24 h). High animal protein intake increases calcium excretion. High calcium intake increases urine calcium. A spot or 24-hour value reflects the diet of the preceding days.
  \item Urine calcium is not a marker of total body calcium or serum calcium. Normal serum calcium with high urine calcium defines hypercalciuria. Low serum calcium with low urine calcium suggests hypocalcemia from calcium deficiency.
  \item Urine calcium must be interpreted with serum calcium and serum creatinine. In hypercalcemia, a low urine calcium suggests FHH. A high urine calcium suggests primary hyperparathyroidism.
  \item Medications alter urine calcium. Thiazide diuretics decrease it. Loop diuretics, glucocorticoids, and excessive vitamin D increase it.
  \item Urine calcium varies with time of day and acid-base status. Metabolic acidosis (renal tubular acidosis type I) increases urine calcium. Alkalosis decreases it.
  \item In children, spot calcium-to-creatinine ratios are used because 24-hour collections are difficult. The normal ratio decreases with age.
\end{itemize}
\section*{Risks}
The test itself is non-invasive and carries no direct risk. The main risk is an invalid result from an incomplete collection. This can lead to unnecessary repeat testing or erroneous management.

\section*{What Happens After the Test}
The 24-hour urine results are available within 1--3 days. The provider interprets urine calcium with serum calcium, serum creatinine, urine creatinine, and the stone type (if known). If hypercalciuria is confirmed, treatment may include dietary modification and thiazide diuretics. Dietary modification includes increasing fluid intake to >2.5 L/day. It includes a normal calcium intake of 1000--1200 mg/day. It includes sodium restriction to <2.3 g/day and limited animal protein. Thiazide diuretics (hydrochlorothiazide 25--50 mg/day or chlorthalidone) reduce urinary calcium. The test is repeated after 3--6 months to assess the response.

\section*{Understanding Results}

\subsection*{Below Normal}
\begin{itemize}
  \item \textbf{Low 24-hour urine calcium (<100 mg/24 h) with hypercalcemia:} This is highly suggestive of familial hypocalciuric hypercalcemia (FHH). This is particularly true if the calcium-to-creatinine clearance ratio is <0.01. FHH is a benign autosomal dominant condition caused by an inactivating mutation of the calcium-sensing receptor (CaSR). It does not require parathyroid surgery.
  \item \textbf{Low urine calcium with normal serum calcium:} This reflects low calcium intake, vitamin D deficiency, or malabsorption. It may coexist with secondary hyperparathyroidism (normal calcium with elevated PTH and low urine calcium).
  \item \textbf{Low urine calcium with hypocalcemia:} This is consistent with hypoparathyroidism, vitamin D deficiency, or calcium deficiency. The kidney appropriately conserves calcium.
  \item \textbf{Low urine calcium in a stone former:} This suggests non-hypercalciuric stone disease (for example, hyperoxaluria, hypocitraturia, uric acid stones) or an incomplete collection. Repeat testing and full stone panel evaluation are indicated.
\end{itemize}

\subsection*{Above Normal}
\begin{itemize}
  \item \textbf{Hypercalciuria (>300 mg/24 h in men, >250 mg/24 h in women, or >4 mg/kg/day):} This is the most common metabolic abnormality in calcium stone formers. Causes include absorptive hypercalciuria, renal hypercalciuria, and resorptive hypercalciuria. Absorptive hypercalciuria is increased intestinal calcium absorption, often diet-related or from vitamin D excess. Renal hypercalciuria is a renal calcium leak that leads to secondary hyperparathyroidism. Resorptive hypercalciuria comes from primary hyperparathyroidism or malignancy. Idiopathic hypercalciuria is also common. It is linked to low bone density in some patients.
  \item \textbf{Hypercalciuria with hypercalcemia and elevated PTH:} This indicates primary hyperparathyroidism. The high urine calcium is caused by an increased filtered load. This combination (hypercalcemia, elevated PTH, normal or high urine calcium) distinguishes primary hyperparathyroidism from FHH. It confirms the need for parathyroidectomy.
  \item \textbf{Hypercalciuria with normal serum calcium:} This indicates idiopathic hypercalciuria or renal calcium leak. This is the classic profile of recurrent calcium oxalate stone formers. Thiazide therapy is highly effective. It reduces urine calcium and stone recurrence.
  \item \textbf{Hypercalciuria with nephrolithiasis and hypocitraturia:} This is a combined metabolic abnormality (hypercalciuria + low urine citrate). It increases the risk of calcium phosphate and calcium oxalate stones. Treatment includes thiazides plus potassium citrate.
  \item \textbf{Hypercalciuria in a patient on glucocorticoids or loop diuretics:} This is drug-induced hypercalciuria. Corticosteroids increase bone resorption and renal calcium excretion. Loop diuretics inhibit tubular calcium reabsorption. Monitor bone density and consider calcium-sparing therapy.
\end{itemize}

\section*{Normal Results}
\begin{itemize}
  \item \textbf{24-hour urine calcium (adults):} 100--300 mg/24 h (2.5--7.5 mmol/24 h) in men; 100--250 mg/24 h (2.5--6.25 mmol/24 h) in women
  \item \textbf{24-hour urine calcium (children):} 1--4 mg/kg/day
  \item \textbf{Urine calcium-to-creatinine ratio (adults, spot):} <0.20 mg/mg (<220 mg/g or <0.56 mmol/mmol)
  \item \textbf{Calcium-to-creatinine clearance ratio (CCCR):} >0.01 in primary hyperparathyroidism; <0.01 in familial hypocalciuric hypercalcemia
\end{itemize}

\section*{Follow-up Tests}
\begin{itemize}
  \item \textbf{Serum calcium, PTH, phosphate, and 25-OH vitamin D:} These tests interpret urine calcium. In hypercalcemia, PTH distinguishes primary hyperparathyroidism (high/inappropriately normal PTH) from other causes (suppressed PTH). In stone formers, these tests identify resorptive hypercalciuria.
  \item \textbf{Full 24-hour urine stone panel:} This includes calcium, oxalate, citrate, uric acid, sodium, creatinine, and volume (2.5--3 L/day target). Hyperoxaluria, hypocitraturia, hyperuricosuria, and low urine volume each require specific management.
  \item \textbf{Stone analysis:} The passed stone is analyzed by infrared spectroscopy or X-ray diffraction. This identifies the composition (calcium oxalate, calcium phosphate, uric acid, cystine, struvite). This guides targeted therapy.
  \item \textbf{CT KUB or renal ultrasound:} These assess stone burden, location, and hydronephrosis. Repeat imaging monitors for new stone formation.
  \item \textbf{Bone density (DEXA) scan:} This is used in patients with hypercalciuria or primary hyperparathyroidism. Long-standing hypercalciuria and hyperparathyroidism cause bone loss.
  \item \textbf{Repeat 24-hour urine collection 3--6 months after starting therapy:} This verifies that thiazide or dietary therapy has reduced urine calcium. It guides dose adjustment. Additional collections may be needed if stones recur.
\end{itemize}

\section*{References}
\bibliographystyle{unsrtnat}
\bibliography{references}

\clearpage
\section*{Regulatory Status and Authorization Date}
This chapter describes a test category, not a specific commercial product. FDA, CDSCO, EU IVDR, and MHRA approval or registration dates therefore cannot be assigned without the assay manufacturer, product name, instrument, and intended use. For laboratory-developed tests, an individual agency approval date may not exist. See the Preface for the interpretation of regulatory status.

\clearpage
